\documentclass[11pt]{article}

\usepackage[margin=1in]{geometry}
\usepackage[utf8]{inputenc}
\usepackage[T1]{fontenc}
\usepackage{booktabs}
\usepackage{graphicx}
\usepackage{url}
\usepackage[colorlinks=true, linkcolor=blue, citecolor=blue, urlcolor=blue]{hyperref}
\usepackage{natbib}

\title{Interventions per 1,000 Robot-Hours:\\
Toward Standardized Reliability Evidence for Deployed Robot Fleets}

\author{
  Manuel Bulos\\
  \texttt{manuelbulos@gmail.com}
}

\date{July 2026}

\begin{document}

\maketitle

\begin{abstract}
Learning-based robots are entering commercial deployment, yet the robotics
industry lacks a standardized way to measure, record, or report their field
reliability. Intervention logs---records of humans pausing, resetting,
teleoperating, or physically assisting a robot---are proprietary,
unstructured, and mutually incompatible across vendors. As a consequence, no
one can answer the question that operators, insurers, and certification
bodies are beginning to ask: \emph{how many human interventions does a robot
fleet require per 1,000 robot-hours, and why?} This position paper argues
that the robotics industry needs the equivalent of aviation's flight data
recorder and the autonomous-vehicle industry's disengagement report, before
regulation imposes one. We propose (i) a minimal set of fleet reliability
metrics centered on the intervention rate per 1,000 robot-hours; (ii) an
open recording schema that captures robot telemetry, policy actions, and
human interventions on a shared clock; and (iii) an intervention taxonomy
with root-cause labels. We release a reference implementation, including a
simulated deployment that demonstrates cross-version regression detection
from intervention telemetry alone. We close with a discussion of the
incentive, privacy, and standardization challenges, and argue that shared
failure evidence is a public good the field should build deliberately rather
than reactively.
\end{abstract}

\section{Introduction}

Commercial deployments of learning-based robots crossed a threshold in
2025--2026: humanoid and manipulation systems now run production shifts in
automotive plants and logistics facilities \citep{figure2026bmw,
agility2025gxo}. The public retrospectives of these pilots are notable both
for what they report and for what they omit. Figure's eleven-month
deployment at BMW Spartanburg reports more than 1,250 hours of runtime and a
task success rate near 98.7\%---which, at production cadence, implies a
failure roughly every 77 cycles, and on the order of hundreds of manual
interventions over the program \citep{figure2026bmw, ifactory2026bmw}. The
deployment's stated goal of zero interventions per shift is not reported as
achieved. Independent reviews of factory-floor humanoid programs find that
intervention logs are treated as proprietary, that no standardized logging
format exists, and that basic reliability questions---mean time between
failures under production conditions, interventions per unit of operating
time, root-cause distributions---cannot be answered from public information
for any commercial robot fleet \citep{karp2026reality}.

This is not a temporary reporting gap; it is a structural one. Robot vendors
have every short-term incentive to keep failure data private, no shared
schema in which to record it even if they wished to publish, and no agreed
metric that would make two fleets comparable. The result is an industry
about to scale physical systems into human environments with less
standardized reliability evidence than aviation had in the 1960s or the
autonomous-vehicle (AV) industry had in 2015.

Two adjacent industries suggest what comes next. In aviation, flight data
recorders became mandatory equipment decades ago, and standardized incident
evidence is credited as a foundation of the field's safety culture
\citep{icao2016annex6}. In the AV industry, California's requirement that
permit holders report \emph{disengagements}---events where a safety driver
takes control---created, whatever its flaws, the first public,
cross-vendor, longitudinal reliability dataset in autonomy
\citep{dmv2024disengagement}. Robotics has no equivalent of either
instrument, while its regulatory context is moving quickly: ISO~25785-1,
the first international safety standard specifically addressing industrial
mobile robots with actively controlled stability (a category that includes
humanoids), is in draft with publication expected around 2027
\citep{iso25785}, and insurers underwriting robot deployments already
demand documented risk assessments and audit trails as conditions of
coverage \citep{humanoidliability2026}.

This paper takes a position: \textbf{the robotics industry should adopt a
standardized, vendor-neutral format for recording and reporting field
reliability evidence---now, while the installed base is small, rather than
after regulation or a high-profile incident forces a worse version of the
same requirement.} We make three concrete proposals and release a reference
implementation.

\section{The evidence gap}

\subsection{What deployments report today}

Public reliability reporting from commercial robot deployments consists
largely of vendor-selected aggregates: cumulative units shipped, cumulative
tasks completed, and task success rates over unspecified conditions
\citep{karp2026reality}. Success rates without exposure denominators
(operating hours, cycles attempted) and without intervention accounting are
not reliability evidence; a 98\% success rate is compatible with both a
robust system and one requiring constant supervision, depending on cycle
time and recovery cost. The metrics that industrial buyers actually use to
evaluate equipment---mean time between failures (MTBF), interventions per
unit time, service burden, uptime---are either not recorded in comparable
form or not disclosed \citep{newmarketpitch2026}.

\subsection{Why the gap persists}

Three forces sustain the gap. \emph{Incentives:} in a competitive funding
environment, no vendor benefits from unilaterally publishing failure rates.
\emph{Infrastructure:} even willing vendors have no shared schema; each
team's logging is an internal engineering artifact, typically retrofitted
onto streaming pipelines that were designed for delivery rather than
faithful recording \citep{roboticstomorrow2026webrtc}. \emph{Demand
immaturity:} until recently, no external party had standing to demand
better. This is changing: insurers now audit deployments and void coverage
on mismatches between the assessed and actual application
\citep{humanoidliability2026}, and conformity assessment against
ISO~25785-1 will require operational records that today's logging practices
cannot produce \citep{iso25785}.

\subsection{The cost of the gap}

The absence of standardized failure evidence is not only a transparency
problem. It slows the field's own progress. Datasets collected by skilled
operators who rarely fail produce policies that never learn recovery
behavior; deliberately captured failure and recovery episodes measurably
improve policy robustness \citep{evsint2026teleop}. Real-deployment
failures---with the full telemetry context that preceded them---are among
the scarcest and most valuable training signals in embodied AI, and today
they are discarded or siloed at the moment of capture.

\section{Proposal 1: A minimal metric set}

We propose that deployed robot fleets be characterized by the following
metrics, all computable from the schema in Section~\ref{sec:schema}.

\begin{table}[h]
\centering
\small
\begin{tabular}{@{}lll@{}}
\toprule
\textbf{Metric} & \textbf{Definition} & \textbf{Unit} \\
\midrule
Intervention rate & Human interventions per 1,000 robot-hours & $\mathrm{I/kh}$ \\
MTBF & Robot-hours between interventions & hours \\
Autonomous success rate & Episodes completed without intervention & fraction \\
Recovered rate & Episodes completed only via human recovery & fraction \\
Mean recovery time & Mean duration of intervention events & seconds \\
Operator burden & Total operator time per 1,000 robot-hours & hours/kh \\
\bottomrule
\end{tabular}
\caption{Proposed fleet reliability metrics. Robot-hours are wall-clock
hours in which the robot was tasked, including time under intervention.}
\label{tab:metrics}
\end{table}

The central quantity is the \textbf{intervention rate per 1,000
robot-hours} ($\mathrm{I/kh}$). We argue for it, rather than task success
rate, as the headline reliability metric for three reasons. First, it has
an exposure denominator, making fleets of different sizes and duty cycles
comparable. Second, it is behaviorally meaningful to operators: it
translates directly into supervision staffing. Third, it has precedent:
it is the robotics analogue of the AV industry's miles-per-disengagement,
a metric that, despite legitimate criticisms of gameability
\citep{dmv2024disengagement}, successfully created public, longitudinal,
cross-vendor comparability where none existed.

An episode that requires human recovery must be classified as
\emph{recovered}, not \emph{successful}. Conflating the two is the single
most common distortion in current vendor reporting.

\section{Proposal 2: A shared-clock recording schema}
\label{sec:schema}

Reliability metrics are only as trustworthy as the records beneath them. We
propose a minimal schema of four append-only streams, all timestamped
against a single monotonic session clock:

\begin{enumerate}
  \item \textbf{Frames}: proprioceptive state at a fixed rate---joint
  positions and velocities, end-effector pose, gripper force and aperture,
  and references to externally stored camera frames.
  \item \textbf{Actions}: policy outputs, each tagged with the model
  version that produced it, the task phase, and optionally the policy's
  own confidence.
  \item \textbf{Interventions}: human take-over events with start and end
  times, intervention type, operator identity, task phase, and a
  root-cause label (Section~\ref{sec:taxonomy}).
  \item \textbf{Episodes}: task attempts with outcome
  (\emph{success}, \emph{recovered}, \emph{failure}, \emph{aborted}) and
  the model version under which they ran.
\end{enumerate}

Two design decisions matter more than the specific fields. First, the
\emph{shared clock}: interventions must be temporally alignable with the
telemetry that preceded them, or root-cause analysis degenerates into
anecdote. Retrofitted logging pipelines routinely fail this requirement
\citep{roboticstomorrow2026webrtc}. Second, \emph{model-version tagging on
every action}: learning-based systems change behavior with every policy
update, and reliability claims that do not condition on model version are
not reproducible. Version tagging is what turns a reliability log into a
regression-detection instrument (Section~\ref{sec:implementation}).

The schema is deliberately storage-format agnostic; our reference
implementation uses JSONL for auditability, and the streams map directly
onto existing robot-learning dataset conventions such as RLDS and LeRobot
\citep{ramos2021rlds, cadene2024lerobot}, so that reliability recording and
training-data capture can be the same act.

\section{Proposal 3: An intervention taxonomy}
\label{sec:taxonomy}

Cross-fleet comparability requires that ``a human stepped in'' mean the
same thing everywhere. We propose a two-level labeling: a closed set of
intervention \emph{types}, and an open, evolving set of root-cause labels.

\begin{table}[h]
\centering
\small
\begin{tabular}{@{}lp{9.5cm}@{}}
\toprule
\textbf{Type} & \textbf{Definition} \\
\midrule
\texttt{pause}    & Operator halted the robot; no state was changed. \\
\texttt{reset}    & Task cycle was restarted from a defined initial state. \\
\texttt{teleop}   & Operator took manual control of the robot to complete
                    or correct the task. \\
\texttt{physical} & Operator physically altered the scene or the robot
                    (cleared a jam, repositioned an object). \\
\texttt{estop}    & Emergency stop, triggered by a human or a safety system. \\
\bottomrule
\end{tabular}
\caption{Proposed closed set of intervention types.}
\label{tab:taxonomy}
\end{table}

Root-cause labels (e.g., \texttt{grasp\_slip}, \texttt{perception\_miss},
\texttt{joint\_fault}, \texttt{collision}) will inevitably be
operator-dependent and incomplete. For this reason the schema records the
telemetry window preceding every intervention, so that unsupervised
clustering over failure signatures can recover root-cause structure even
where labels are missing or inconsistent---an approach we demonstrate in
the reference implementation.

\section{Reference implementation}
\label{sec:implementation}

We release \texttt{robot-blackbox}, an open-source reference implementation
of the schema and metrics.\footnote{\url{https://github.com/manuelbulos/robot-blackbox}}
It comprises a recording SDK whose API wraps an arbitrary control loop; the
metric computations of Table~\ref{tab:metrics}; failure clustering over
pre-intervention telemetry windows; cross-version regression detection; and
an audit-oriented report generator.

To exercise the pipeline end-to-end without hardware, the implementation
includes a simulated 6-DoF manipulation deployment with injectable failure
modes whose telemetry signatures (gripper force collapse on slips, velocity
variance on faults) are modeled on failure classes reported from real
pilots. In a demonstration scenario, a fleet runs 200 cycles on a baseline
policy version and 200 cycles on an updated version with a deliberately
regressed grasping module. From intervention telemetry alone, the pipeline
reports the update's intervention rate as $1.8\times$ the baseline's
(8{,}766 vs.\ 4{,}782 per 1,000 robot-hours in the simulated fleet) and
localizes the regression to the grasp phase via clustering. We emphasize
that these figures are simulated and illustrate the \emph{instrument},
not any real fleet; producing the first real cross-vendor baselines is
precisely the point of standardizing the format.

\section{Discussion}

\textbf{Incentives.} No vendor will publish reliability evidence
unilaterally, and this proposal does not ask them to. The near-term
adopters are the parties who already need the records privately: operators
who must staff supervision, insurers who must price deployments
\citep{humanoidliability2026}, and vendors who must diagnose their own
regressions. Public benchmarks can emerge later from anonymized
aggregation, as they did in the AV industry only under regulatory
compulsion \citep{dmv2024disengagement}. The window to shape the format
voluntarily is before ISO~25785-1 conformity assessment and insurance
requirements harden around ad-hoc alternatives \citep{iso25785}.

\textbf{Privacy and competition.} Intervention records encode competitive
information. The schema therefore separates the \emph{format} (open) from
the \emph{data} (owned by the deploying parties), and aggregation for
benchmarking requires explicit anonymization of task, site, and embodiment
identifiers. The aviation precedent is instructive: standardized recorders
coexist with confidential investigation regimes.

\textbf{Gaming.} Any headline metric invites optimization against it; AV
disengagement counts were criticized on exactly this ground. Two properties
of the proposed schema mitigate this: interventions are recorded with
surrounding telemetry (making deletion detectable as coverage gaps on the
shared clock), and the episode outcome taxonomy makes ``success via human''
unrepresentable as autonomous success.

\textbf{Failure data as a training asset.} Because the schema is compatible
with robot-learning dataset formats \citep{ramos2021rlds, cadene2024lerobot},
every recorded intervention is simultaneously a labeled failure-and-recovery
episode---the training signal that demonstration datasets systematically
lack \citep{evsint2026teleop}. Reliability evidence and recovery-learning
data are the same bytes; the industry currently throws both away.

\section{Conclusion}

Robotics is repeating a pattern aviation and autonomous driving have
already lived through: capability arrives before evidence infrastructure,
and evidence infrastructure arrives only when demanded from outside. The
field has a short window in which the instrument can be designed by the
people who understand robots, rather than imposed by those who do not. A
shared clock, four streams, five intervention types, and one honest
metric---interventions per 1,000 robot-hours---would be enough to start.

\bibliographystyle{plainnat}
\bibliography{references}

\end{document}
